\chapter{Diseño e implementación}
\label{chap:diseno}
En este capítulo se describe la arquitectura del sistema, las estructuras de datos, los algoritmos principales y las decisiones de implementación que sustentan el proyecto, desde la ingesta de datos hasta la simulación y la generación de salidas.
\section{Obtención y preparación de los datos}\label{sec:datos}

Para realizar la simulación fue necesario disponer de varios subconjuntos de datos procedentes del catálogo \textit{Gaia DR3} y del catálogo complementario de distancias \texttt{GeDR3Dist}.
El proceso completo comprendió la descarga, reducción y combinación de los distintos catálogos, con el objetivo de obtener un conjunto coherente de parámetros espaciales, fotométricos y físicos.

\subsection*{Descarga de los catálogos principales}

Los catálogos \texttt{gaia\_source} y \texttt{astrophysical\_parameters} se descargaron desde el \textit{Gaia Archive}~\cite{gaia_archive} mediante el comando \texttt{wget}, automatizando la obtención de archivos comprimidos en formato \texttt{.csv.gz}:

\begin{verbatim}
wget -r -np -nd -A "*.csv.gz" "URL_archivos"
\end{verbatim}

Cada archivo contiene un subconjunto del catálogo particionado por rangos de \texttt{source\_id}, lo que facilita la descarga y el procesamiento distribuido.
El tamaño total de la descarga es de aproximadamente 757~GB para la tabla \texttt{gaia\_source} y 258~GB para \texttt{astrophysical\_parameters}, sumando un total de 3\,386 archivos en cada caso.
\subsection*{Descarga del catálogo de distancias}

Las estimaciones de distancia derivadas del paralaje se obtuvieron a partir del catálogo \texttt{I/352/GeDR3Dist}, disponible a través del servicio \textit{TAPVizieR}~\cite{tap_vizier}.
Dado el volumen de datos (más de mil millones de fuentes), se diseñó un script en \texttt{bash}~\ref{lst:script_descargar_distances} que divide el rango de \texttt{source\_id} en bloques y realiza consultas ADQL (\textit{Astronomical Data Query Language}) sucesivas al servidor.

De esta forma se obtuvieron diversos archivos \texttt{.csv} con los campos \texttt{Source} y \texttt{rgeo}, que posteriormente se integraron en el conjunto de datos principal.
\subsection*{Reducción de columnas}

Con el fin de optimizar el almacenamiento y acelerar las operaciones de lectura, los archivos originales se redujeron manteniendo únicamente las columnas relevantes para el análisis:
\texttt{source\_id}, \texttt{ra}, \texttt{dec}, \texttt{pmra}, \texttt{pmdec}, \texttt{radial\_velocity}, \texttt{phot\_g\_mean\_mag} y \texttt{bp\_rp}.
El filtrado se realizó mediante herramientas de línea de comandos (\texttt{awk}) y utilizando \texttt{parallel} para acelerar el proceso, generando versiones compactas de los catálogos~\ref{lst:script_reducir_astro}.
\subsection*{Combinación de catálogos}

Finalmente, los catálogos reducidos de \texttt{gaia\_source}, \texttt{astrophysical\_parameters} y \texttt{GeDR3Dist} se combinaron utilizando un script en \texttt{Python}~\ref{lst:script_combinar} con las librerías \texttt{pandas} y \texttt{duckdb}.
Primero se creó una base de datos temporal con las distancias debido al diferente número de archivos por tabla y el orden de IDs, y posteriormente se fusionaron los catálogos mediante uniones internas sobre el campo \texttt{source\_id}, aprovechando el procesamiento en paralelo para reducir el tiempo de ejecución:

\begin{verbatim}
merged = pd.merge(main_df, astro_df, on='source_id')
merged = pd.merge(merged, distance_df, on='source_id')
\end{verbatim}

El resultado final fue un conjunto de 3386 archivos optimizados \texttt{merged\_XXXX.csv}, listos para el filtrado y el análisis posterior ocupando 223~GB sin comprimir.

Este procedimiento permitió manejar de forma eficiente el volumen total del catálogo (más de 1\,800 millones de fuentes en \textit{Gaia DR3}), garantizando coherencia y consistencia entre todos los parámetros empleados.
\section{Arquitectura del sistema y flujo de datos}\label{sec:arquitectura}
El sistema implementa un simulador de dinámica galáctica mediante una arquitectura modular que soporta dos implementaciones complementarias: una versión \textbf{CPU}, basada en paralelización mediante OpenMP, y una versión \textbf{GPU}, que aprovecha la capacidad de cómputo masivo de aceleradores NVIDIA mediante CUDA\@.

Ambas versiones comparten la misma interfaz de entrada/salida, las mismas estructuras de datos fundamentales y el mismo pipeline de procesamiento, permitiendo una fácil validación cruzada de resultados y comparación de rendimiento.
La selección entre la versión CPU o GPU se realiza en tiempo de compilación mediante opciones de CMake (\texttt{ENABLE\_CUDA}), facilitando la adaptación a distintas infraestructuras y volumen de datos.


FALTA COMPLETAR

\section{Estructuras de datos principales}\label{sec:estructuras}
El sistema utiliza varias estructuras de datos fundamentales que son compartidas entre la versión CPU y GPU, además de otras especializadas en cada versión.
Estas estructuras están diseñadas para optimizar tanto el acceso a memoria como la compatibilidad con los modelos de paralelización disponibles.

\subsection{Estructura de estrella (\texttt{Star})}\label{subsec:estructura-de-estrella-(texttt{star})}

La estructura central del sistema es \texttt{Star}, que agrupa los parámetros astrométricos, cinemáticos y físicos de cada objeto celeste.
Esta estructura se organiza de forma \textit{structure-of-arrays} (SoA) en lugar de arrays de estructuras, permitiendo optimizar el acceso a memoria en operaciones paralelas y facilitar la vectorización automática en CPU y GPU.

\subsubsection*{Diseño y selección de tipos de datos}

La selección de tipos de datos en cada campo responde a criterios de precisión numérica y eficiencia de memoria:

\begin{itemize}
    \item \textbf{Precisión doble (\texttt{double}, 64 bits)}:
    \begin{itemize}
        \item \texttt{ra}, \texttt{dec}, \texttt{distance}, \texttt{pmra}, \texttt{pmdec}, \texttt{radial\_velocity}: Los datos originales de Gaia DR3 se proporcionan en precisión doble. Mantenerlos así preserva la precisión astrométrica y es necesario para transformaciones de coordenadas que acumulan errores.

        \item \texttt{Cx}, \texttt{Cy}, \texttt{Cz}, \texttt{Vx}, \texttt{Vy}, \texttt{Vz}: Las coordenadas cartesianas heliocéntricas y velocidades se mantienen en doble precisión para garantizar estabilidad numérica durante la integración temporal de órbitas a lo largo de millones de años. Los errores de redondeo en precisión simple acumularían desviaciones significativas en las trayectorias.
    \end{itemize}

    \item \textbf{Precisión simple (\texttt{float}, 32 bits)}:
    \begin{itemize}
        \item \texttt{mass}, \texttt{radius}, \texttt{gravity}, \texttt{mean\_g}, \texttt{color}: Estos parámetros representan propiedades físicas derivadas o fotométricas con precisión inherente limitada. La masa se estima mediante relaciones empíricas con incertidumbre del orden del 10--20\%. El índice de color y magnitud se proporcionan con precisión limitada en el catálogo. Usar precisión simple reduce el consumo de memoria a la mitad en estos campos sin pérdida práctica de información.
    \end{itemize}

    \item \textbf{Identificador de 64 bits sin signo (\texttt{unsigned long})}:
    \begin{itemize}
        \item \texttt{id}: Identificador único de Gaia DR3, necesario para relacionar datos con el catálogo original.
    \end{itemize}
\end{itemize}

\subsubsection*{Consumo de memoria}
Para una estructura con $N$ estrellas y capacidad $C$:
\[
    \text{Memoria} = C \times (8 + 8 \times 6 + 8 \times 3 + 8 \times 3 + 4 \times 5) \text{ bytes} = C \times 152 \text{ bytes}
\]

Esto es:
\begin{itemize}
    \item \texttt{id}: 8 bytes $\times$ 1 = 8 bytes
    \item \texttt{ra}, \texttt{dec}, \texttt{distance}, \texttt{pmra}, \texttt{pmdec}, \texttt{radial\_velocity}: 8 bytes $\times$ 6 = 48 bytes
    \item \texttt{Cx}, \texttt{Cy}, \texttt{Cz}, \texttt{Vx}, \texttt{Vy}, \texttt{Vz}: 8 bytes $\times$ 6 = 48 bytes
    \item \texttt{mean\_g}, \texttt{color}, \texttt{radius}, \texttt{gravity}, \texttt{mass}: 4 bytes $\times$ 5 = 20 bytes
    \item \textbf{Total por estrella: 152 bytes}
\end{itemize}

Para una simulación con 1.100 millones de estrellas utilizadas en el proyecto como se ve en la subsección~\ref{subsec:filtrado}, la memoria requerida sería:

\[
    1.1 \times 10^9 \times 152 \text{ bytes} = 171.3 \text{ GB}
\]
tamaño perfectamente abarcable en los nodos disponibles en el Finisterrae III de 256 GB de RAM~\ref{subsec:especificaciones-tecnicas-de-los-nodos}.

\subsubsection*{Liberación selectiva de datos}

Tras la carga completa de un bloque de estrellas y la realización de los cálculos preliminares (transformación de coordenadas y estimación de masas), ciertos campos astrométricos originales dejan de ser necesarios y son liberados para reducir el consumo de memoria:

\begin{itemize}
    \item \textbf{Datos liberados después de transformación de coordenadas}:
    \begin{itemize}
        \item \texttt{ra}, \texttt{dec}, \texttt{distance}: Tras convertir a coordenadas cartesianas heliocéntricas (\texttt{Cx}, \texttt{Cy}, \texttt{Cz}), estos datos ecuatoriales ya no son necesarios.

        \item \texttt{pmra}, \texttt{pmdec}, \texttt{radial\_velocity}: Tras convertir a velocidades cartesianas (\texttt{Vx}, \texttt{Vy}, \texttt{Vz}), estos datos originales no vuelven a usarse.
    \end{itemize}

    \item \textbf{Datos liberados después de estimación de masas}:
    \begin{itemize}
        \item \texttt{gravity}, \texttt{mean\_g}, \texttt{color}: Tras estimar \texttt{mass} mediante la relación masa-luminosidad que utiliza estos parámetros, pueden ser liberados si no se requieren en resultados finales.
    \end{itemize}

    \item \textbf{Impacto en memoria}:
    Liberando los 6 campos astrométricos originales reduce el consumo en 48 bytes por objeto, correspondiendo a una reducción del 31.6\% respecto al consumo total por estrella.
    Para 1.100 millones de estrellas, esto libera aproximadamente 49.2 GB de memoria.
\end{itemize}

Este diseño permite un balance entre precisión numérica, eficiencia de memoria y adaptabilidad a diferentes infraestructuras y volúmenes de datos, siendo crucial para gestionar simulaciones de escala galáctica en entornos de computación de altas prestaciones.

\subsection{Estructura de arbol CPU(\texttt{Octree})}\label{subsec:estructura-de-arbol-cpu(texttt{octree})}

La estructura \texttt{Octree} implementa el árbol jerárquico utilizado en la versión CPU del simulador para acelerar el cálculo de fuerzas gravitatorias mediante el método de Barnes-Hut.
A diferencia de la estructura \texttt{Star}, que es plana y homogénea, el octree es una estructura (SoA) jerárquica donde cada nodo contiene referencias a hasta 8 nodos hijos.

\subsubsection*{Diseño y campos de datos}

Cada nodo del octree almacena información geométrica, propiedades físicas acumuladas e información topológica:

\begin{itemize}
    \item \textbf{Geometría del volumen (precisión simple, \texttt{float})}:
    \begin{itemize}
        \item \texttt{center\_x}, \texttt{center\_y}, \texttt{center\_z}: Coordenadas del centro del cubo (float, 3 $\times$ 4 = 12 bytes)

        \item \texttt{half\_size}: Semiextensión del cubo en cada dirección (float, 4 bytes)

        \textbf{Justificación de precisión simple}: La geometría del árbol se utiliza únicamente para aplicar el criterio de apertura de Barnes-Hut (comparación de distancias). La precisión simple es suficiente para estas comparaciones relativas sin afectar la validez del criterio de agrupación.
    \end{itemize}

    \item \textbf{Propiedades físicas (masa en \texttt{float}, centro de masas en \texttt{double})}:
    \begin{itemize}
        \item \texttt{mass}: Masa total acumulada del nodo y sus descendientes (float, 4 bytes)

        \item \texttt{com\_x}, \texttt{com\_y}, \texttt{com\_z}: Centro de masas (double, 3 $\times$ 8 = 24 bytes)

        \textbf{Justificación de precisión mixta}: La masa total se almacena en precisión simple ya que es el resultado de sumas acumulativas que se recomputabilizan en cada reconstrucción del árbol.
        El centro de masas, sin embargo, se mantiene en precisión doble porque:
        \begin{enumerate}
            \item Se utiliza directamente en el cálculo de la aceleración gravitatoria.
            \item Se actualiza mediante una operación de media ponderada que puede acumular errores significativos en precisión simple.
            \item La precisión en la localización del centro de masas afecta directamente a la exactitud de la aproximación multipolar.
        \end{enumerate}
    \end{itemize}

    \item \textbf{Información topológica}:
    \begin{itemize}
        \item \texttt{children}: Array de 8 referencias a nodos hijos (array de 8 enteros sin signo de 32 bits, 32 bytes).

        Se utiliza \texttt{unsigned int} en lugar de \texttt{long} por dos razones críticas:
        \begin{enumerate}
            \item \textbf{Ahorro de memoria}: Cada referencia ocupa 4 bytes en lugar de 8, reduciendo en un 50\% el espacio dedicado a punteros.

            \item \textbf{Límite máximo de nodos}: El rango de \texttt{unsigned int} es de 0 a $2^{32} - 1 = 4{,}294{,}967{,}295$, más que suficiente para los nodos esperados en una simulación de 1.100 millones de estrellas.
        \end{enumerate}

        La representación de sin hijo utiliza la constante \texttt{INVALID\_INDEX = UINT32\_MAX} (4.294 mil millones), un valor fuera del rango de nodos válidos.
        Aunque \texttt{unsigned int} no permite directamente valores negativos como -1, el uso de \texttt{UINT32\_MAX} es una estrategia estándar que proporciona semántica equivalente: una comparación \texttt{if (child == INVALID\_INDEX)} determina de forma unívoca si el hijo existe.

        \item \texttt{star\_index}: Índice de la estrella si el nodo es una hoja, o -1 si es un nodo interno (long, 8 bytes).

        Este campo mantiene precisión de 64 bits porque son los mismos valores que los almacenados en Star.
    \end{itemize}
\end{itemize}

\subsubsection*{Consumo de memoria}

Memoria por nodo del octree:

\[
    \text{Memoria por nodo} = 4 \times 4 + 4 + 24 + 32 + 8 = 88 \text{ bytes}
\]

Esto es:
\begin{itemize}
    \item \texttt{center\_x}, \texttt{center\_y}, \texttt{center\_z}, \texttt{half\_size}: 4 bytes $\times$ 4 = 16 bytes
    \item \texttt{mass}: 4 bytes
    \item \texttt{com\_x}, \texttt{com\_y}, \texttt{com\_z}: 8 bytes $\times$ 3 = 24 bytes
    \item \texttt{children[8]}: 4 bytes $\times$ 8 = 32 bytes
    \item \texttt{star\_index}: 8 bytes
    \item \textbf{Total por nodo: 88 bytes}
\end{itemize}

\subsection{Estructura de arbol GPU(\texttt{OctreeGPU})}\label{subsec:estructura-de-arbol-gpu(texttt{octreegpu})}
La estructura \texttt{OctreeGPU} implementa el árbol jerárquico para la versión acelerada por GPU\@.
Aunque mantiene la misma lógica algorítmica de Barnes-Hut que la versión CPU, su organización física es fundamentalmente diferente, optimizada para:
\begin{enumerate}
    \item Maximizar la coalesced memory access en GPU (acceso contiguo a memoria)
    \item Minimizar las divergencias de control entre threads CUDA
    \item Reducir la necesidad de sincronización global entre dispositivos
    \item Facilitar la paralelización masiva de búsquedas en el árbol
\end{enumerate}

En lugar de un único árbol monolítico como en CPU, la versión GPU divide el árbol raíz en 8 subárboles independientes, uno para cada octante del espacio:

\begin{lstlisting}[language=C,label=octree_gpu_struct,caption=Estructura OctreeGPU con partición en octantes]
typedef struct {
    OctreeOctant *octants[8];    // 8 subárboles independientes
    FrontierGPU *frontier[8];    // 8 fronteras multipolo
} OctreeGPU;
\end{lstlisting}

Esta partición ofrece ventajas significativas:

\begin{itemize}
    \item \textbf{Paralelización natural}: Los 8 octantes se construyen en paralelo durante la etapa de inicialización.
    Cada thread construye un octante sin contención con los otros.

    \item \textbf{Eficiencia de memoria GPU}: Cada octante se transfiere a la GPU como un bloque de memoria contiguo, evitando la busqueda de hijos por rama como en CPU.

    \item \textbf{Escalabilidad multi-GPU}: En sistemas con múltiples aceleradores, cada octante puede asignarse a una GPU diferente, distribuye la carga de trabajo de forma natural.

\end{itemize}

A diferencia de la versión CPU donde cada nodo se almacena de forma distribuida en arrays separados (SoA), en GPU cada nodo individual se organiza como una estructura compacta \texttt{OctreeNode} (\textit{Array of Structures}, AoS) que contiene los mismos datos en una única unidad de memoria.
Esto se debe a que en GPU el costo de divergencia de control al acceder a múltiples arrays dispersos es mayor que el costo de leer una estructura compacta.
Cada thread de CUDA que desciende por el árbol accede a un único nodo a la vez; tener todos los datos del nodo juntos en memoria maximiza la localidad espacial, mejorando la tasa de coalesced memory access.

Debido a la baja eficiencia de funciones recursivas para recorrer el árbol en GPU, la arquitectura CUDA requiere un enfoque iterativo que evite el uso de stack implícito.
El campo añadido a la estructura \texttt{next} implementa un enlace al siguiente nodo relevante en orden de traversal horizontal, convirtiendo la búsqueda recursiva del árbol en una iteración directa sobre una lista encadenada de nodos no-hoja.
\begin{lstlisting}[language=C,label={lst:octree_node_struct},caption=Estructura OctreeNode]
typedef struct {
    float center_x, center_y, center_z;
    float half_size;
    float mass;
    double com_x, com_y, com_z;
    unsigned int children[8];
    long star_index;
    unsigned int next;
} OctreeNode;
\end{lstlisting}

Debido a la subdivisión por octantes, es necesario representar una porción minimal del resto de subárboles, la frontera multipolo cumple este propósito fundamental.
La frontera multipolo \texttt{FrontierGPU} es la estructura de datos central que representa la aproximación del campo gravitatorio generado por el resto del árbol.
En lugar de enumerar explícitamente todos los nodos que contribuyen a la aceleración de una partícula (lo cual generaría divergencia masiva entre threads), se mantiene una lista compacta de nodos que conforman la frontera válida según el criterio de apertura de Barnes-Hut.

\section{Diseño algorítmico y paralelización}\label{sec:diseno-algoritmico-y-paralelizacion}





